\section{Anécdotas, teorías de la conspiración y jugadores subestimados}

Esta sección reúne percepciones e historias dispersas pero muy valiosas de la entrevista, que no pertenecen a un solo tema, pero cada una revela algún aspecto de la industria de la IA.

\subsection{Teorías de la conspiración vs. la verdad}

Las teorías de la conspiración en torno a DeepSeek se propagaron como fuego en las redes sociales. Dylan y Nathan las desmontan una por una:

\textbf{«El gobierno chino está subsidiando a DeepSeek»}: poco probable. DeepSeek está financiada por High-Flyer (幻方量化, un fondo de cobertura), no por una empresa vinculada al gobierno. Liang Wenfeng (梁文锋) sí se reunió después con dirigentes del gobierno, pero eso fue \textbf{después} del éxito de DeepSeek: el gobierno se sumó a la ola, no la financió por detrás.

\textbf{«Vendieron en corto acciones de NVIDIA antes del lanzamiento»}: también poco probable. V3 se lanzó el 26 de diciembre: ¿quién elegiría el primer día después de Navidad para coordinar una venta en corto? \textit{«Creo que simplemente estaban apurados por terminar; a quién le importa la Navidad, con tal de sacarlo antes del Año Nuevo chino.»} (Nathan)

\textbf{«Solo costó 6 millones de dólares entrenarlo»}: es una lectura muy equivocada del artículo. Los \$5.576M son el costo en horas de GPU de \textbf{una sola corrida} de preentrenamiento de V3. No incluye los experimentos previos (posiblemente decenas de ablaciones a menor escala), el entrenamiento posterior, el entrenamiento de R1 ni el servicio de inferencia. La inversión total real podría ser de 10 a 50 veces esa cifra.

\subsection{El subestimado Gemini Flash Thinking}

En las pruebas personales de Lex, Gemini Flash Thinking de Google podría ser más barato que R1 y no menos capaz, pero casi nadie habla de él. Las razones posibles son:
\begin{itemize}
    \item El marketing de Google está muy por debajo, en capacidad de llamar la atención, de la estrategia de código abierto de DeepSeek
    \item Flash Thinking podría usar un método distinto (superponer razonamiento sobre una arquitectura existente, en vez de un entrenamiento de razonamiento dedicado)
    \item La narrativa de DeepSeek como «caballo negro chino» tiene más capacidad de difusión

\end{itemize}

Es una lección sobre «tecnología vs. narrativa»: la mejor tecnología no necesariamente se lleva la mayor atención.

\subsection{PyTorch PowerPlant.NoBlowUp}

Al entrenar clústeres grandes, las fluctuaciones de potencia de las GPU pueden destruir el equipo eléctrico. Fase de cómputo con alta potencia (todas las GPU calculando a toda capacidad) $\to$ fase de comunicación con baja potencia (las GPU esperando el intercambio de datos). Esta carga pulsante es muy peligrosa para los transformadores y el equipo de distribución eléctrica.

Un ingeniero de Meta escribió un operador de PyTorch para resolver este problema: durante los periodos ociosos de las GPU, hacerlas calcular números sin sentido (multiplicaciones falsas), únicamente para mantener una potencia estable. Este código se hizo público por accidente: un operador llamado \texttt{PowerPlant.no\_blowup = 1}.

\textit{«Es muy fácil que termines volando algo en pedazos.»} (Dylan) Esta historia ilustra perfectamente el reto de ingeniería «aburrido pero crítico» del entrenamiento de IA: el problema no es un avance algorítmico, sino cómo evitar que tu transformador eléctrico explote.

\subsection{El dilema de seguridad de Anthropic}

Se reporta que Anthropic tiene un modelo de razonamiento mejor que O3, pero no lo publica por consideraciones de seguridad. La cadena de pensamiento de R1 sí puede resultar inquietante: cambia entre chino e inglés, aparecen fragmentos que parecen galimatías y luego, de pronto, da la respuesta correcta.

El lanzamiento agresivo de DeepSeek redujo el estándar de seguridad de todos, de forma similar a como la presión del programa espacial soviético en la Guerra Fría afectó a la NASA estadounidense. \textit{«DeepSeek ships. That's one of their big advantages.»} (Nathan) El lanzamiento rápido en sí mismo es un arma competitiva.

\subsection{La profecía de Sam Altman}

Se cita una frase de Sam Altman muy repetida: \textit{«La persuasión sobrehumana llegará antes que la inteligencia sobrehumana.»} Esto tiene implicaciones profundas para la seguridad de la IA: antes de que la IA pueda «pensar de verdad», es posible que ya pueda «persuadir de verdad». Los chatbots de Character AI ya influyen en el estado de ánimo de usuarios jóvenes; ¿qué pasaría si fuera una manipulación cultural o política deliberada?

También llama la atención una declaración pública de Zuckerberg en una llamada de resultados financieros: \textit{«Para la ventaja de nuestro país, es importante que el estándar de código abierto sea estadounidense.»} El código abierto no es solo una estrategia técnica, sino también una herramienta geopolítica.

\subsection{Resumen de la sección}

La industria de la IA está llena de detalles fascinantes: desde el reto de ingeniería de «no volar el transformador» hasta la lectura equivocada de los «6 millones de dólares», pasando por el subestimado Gemini y el código abierto como arma geopolítica. Estas anécdotas revelan una realidad de la industria más rica, más caótica y más humana que la de los artículos técnicos.
